% Generated draft: requires mathematical and source review.
\documentclass[11pt]{article}
\usepackage[margin=25mm]{geometry}
\usepackage{fontspec}
\setmainfont{Latin Modern Roman}
\usepackage{amsmath,amssymb,mathtools,unicode-math}
\setmathfont{Latin Modern Math}
\usepackage{enumitem,needspace,xurl,hyperref}
\setlength{\parindent}{0pt}
\setlength{\parskip}{5pt}
\setlength{\emergencystretch}{3em}
\newcommand{\N}{\mathbb{N}}
\newcommand{\Z}{\mathbb{Z}}
\newcommand{\Q}{\mathbb{Q}}
\newcommand{\R}{\mathbb{R}}
\newcommand{\C}{\mathbb{C}}
\newcommand{\F}{\mathbb{F}}
\newcommand{\A}{\mathbb{A}}
\newcommand{\PP}{\mathbb{P}}
\newcommand{\E}{\mathbb{E}}
\newcommand{\eps}{\varepsilon}
\newcommand{\dd}{\,\mathrm{d}}
\begin{document}
\section*{UnsolvedMath ID: 30006003}
\textbf{Source problem code:} catalogue entry
\subsection{Definitions and mathematical statement}
The supplied record does not define \(c_{\mathrm{BE}}\), its function space, or its distance convention. The following is a provisional formulation using the conventional real-valued classical Sobolev setting; these choices require confirmation against the cited report.

For an integer \(d\geq 3\), set \(2^*=2d/(d-2)\). Let \(\dot H^1(\R^d;\R)\) be the completion of \(C_c^\infty(\R^d;\R)\) in the norm \(u\mapsto\|\nabla u\|_{L^2}\), represented by functions in \(L^{2^*}(\R^d)\). Define the sharp Sobolev constant by
\[
S_d=\inf_{0\ne u\in\dot H^1(\R^d;\R)}
\frac{\|\nabla u\|_{L^2}^2}{\|u\|_{L^{2^*}}^2}.
\]
In this conventional setting, let
\[
\mathcal M_d=\left\{x\longmapsto a\left(\frac{\lambda}{\lambda^2+|x-y|^2}\right)^{(d-2)/2}:a\in\R,\ \lambda>0,\ y\in\R^d\right\},
\qquad
\delta_d(u)=\inf_{v\in\mathcal M_d}\|\nabla(u-v)\|_{L^2}.
\]
Thus \(\mathcal M_d\) is the conventional family of Sobolev equality functions, including zero. Define
\[
Q_d(u)=\frac{\|\nabla u\|_{L^2}^2-S_d\|u\|_{L^{2^*}}^2}{\delta_d(u)^2}
\quad\text{when }\delta_d(u)>0,
\qquad
c_{\mathrm{BE}}(d)=\inf_{\substack{u\in\dot H^1(\R^d;\R)\\\delta_d(u)>0}}Q_d(u).
\]
An optimizer of the stability constant is a function in
\[
\mathcal O_d=\{u\in\dot H^1(\R^d;\R):\delta_d(u)>0,\ Q_d(u)=c_{\mathrm{BE}}(d)\}.
\]
The positive-distance condition distinguishes these optimizers from the Sobolev equality functions.

\textbf{Problem, under this interpretation.} For each integer \(d\geq3\), characterize all elements of \(\mathcal O_d\). In particular, determine whether every \(u\in\mathcal O_d\) is radial about some point, meaning that there exist \(y\in\R^d\) and a measurable function \(g:[0,\infty)\to\R\) such that \(u(x)=g(|x-y|)\) almost everywhere. Also determine whether there exists \(u\in\mathcal O_d\) for which both \(\{x:u(x)>0\}\) and \(\{x:u(x)<0\}\) have positive Lebesgue measure. Any dependence of these answers on \(d\) is part of the characterization.

The dimensional range above is suggested by the supplied literature assessment, rather than explicitly specified in the extracted question. The interpretation of radiality as radiality about some point is also provisional. The supplied dated curation reports existence of an optimizer in every classical dimension \(d\geq3\), while describing characterization and qualitative properties as unresolved; this draft does not independently verify those claims.
\subsection{Short English statement}
Characterize the functions attaining the sharp Bianchi--Egnell stability constant, including whether all are radial about some point and whether any change sign. The displayed formulation uses a provisional classical real-valued Sobolev setting because the supplied extract omits the defining quotient and function space. The supplied partially solved status and existence assessment require independent verification.
\subsection{Sources}
\begin{itemize}
\item Local UnsolvedMath registry, permanent ID \(30006003\), title \emph{Bianchi-Egnell Optimizers}, source problem code \texttt{catalogue entry}, in the Oberwolfach Reports Open Problems collection. The record gives proposed year 2024 and names no proposer.
\item Supplied report citation: \emph{Calculus of Variations (2025)}, Oberwolfach Reports, DOI \texttt{10.4171/owr/2024/37}, pp.\ 2113--2178. \url{https://doi.org/10.4171/owr/2024/37}. The background describes Oberwolfach Reports as documenting mathematical workshops and their open-problem discussions.
\item Supplied literature reference, labelled ``Konig, existence of a minimizer'': \url{https://ems.press/content/serial-article-files/49782}.
\item Supplied literature reference, labelled ``Konig, sharp-constant bounds'': \url{https://arxiv.org/abs/2210.08482}.
\item Literature assessment checked at \texttt{2026-08-22}: supplied dated curation, not independently verified. It labels the problem partially solved and reports that a minimizer attaining \(c_{\mathrm{BE}}\) exists in every classical dimension \(d\geq3\), while characterization and qualitative properties of all sharp-constant optimizers remain open. The background explicitly requests independent verification of time-sensitive status and cited claims.
\item The registry labels the statement \texttt{corrected\_verified} and says it was reconstructed and checked against the official Oberwolfach report. These are supplied provenance claims, not verification performed for this draft. The retained original statement asks for characterization and basic qualitative properties; the reconstructed statement specifically adds radiality and sign change.
\end{itemize}
\end{document}
